% Non-Diagonal Kinetic Matrix: Alternative Approaches
% Source: docs/kinetic_matrix_alternatives.md (migrated to LaTeX)
% Uses macros from preamble.tex

\section{Non-Diagonal Kinetic Matrix: Alternative Approaches}
\label{sec:kinetic-matrix}

When the kinetic matrix $K_{ij} = \partial \pi_i / \partial \dot{q}_j$ is non-diagonal (e.g., linearized gravity, massive gravity), the standard Hamiltonian split $\mathrm{d}q/\mathrm{d}t = \pi$ fails. This document records all investigated approaches and their tradeoffs.


\subsection{Current Implementation: Symbolic $K^{-1}$ Inversion (Wolfram)}
\label{sec:kinetic-current}

\textbf{Status:} Implemented and working (commit \texttt{906b250}, hardened in \texttt{cb7c12f}).

Compute $K$ symbolically in Wolfram, invert via \texttt{Inverse[K]}, emit $K^{-1}_{ij}$ coefficients to JSON. Field rates become:
\begin{equation}
  \frac{\mathrm{d}q_i}{\mathrm{d}t} = \sum_j K^{-1}_{ij}\, (\pi_j - S_j)
  \label{eq:kinv-field-rates}
\end{equation}
where $S_j$ is the spatial part of $\pi_j$ (momenta evaluated with all velocities $= 0$).

\textbf{Pros:}
\begin{itemize}
  \item Gauge-independent --- works for any theory without gauge fixing
  \item Exact symbolic inversion --- no numerical errors
  \item Pre-computed at derive-time --- zero runtime cost
  \item Handles constant and parameter-dependent $K$ (via \texttt{coefficient\_symbolic})
\end{itemize}

\textbf{Cons:}
\begin{itemize}
  \item $K^{-1}$ entries can be large rational expressions for big systems
  \item Requires $\det(K) \neq 0$ (throws if singular, indicating constraint structure)
  \item \emph{Position-dependent $K$ (from curved metrics) would need grid evaluation at runtime}
\end{itemize}

\textbf{Works for:} Proca ($2\times 2$ diagonal), massive gravity ($4\times 4$), gravitational waves ($7\times 7$ dynamical block), elasticity ($2\times 2$ diagonal with symbolic $\rho$).


\subsection{Alternative 1: Mass Matrix ODE Solvers (Julia / PETSc / SUNDIALS)}
\label{sec:kinetic-alt1}

\textbf{Status:} Investigated, not implemented. Viable future option.

Instead of inverting $K$, reformulate as a mass-matrix ODE:
\begin{equation}
  M \cdot \frac{\mathrm{d}y}{\mathrm{d}t} = f(y, t)
  \label{eq:mass-matrix-ode}
\end{equation}
where $M$ encodes the kinetic coupling. Specialised solvers handle $M$ directly.

\subsubsection{Julia DifferentialEquations.jl}
\label{sec:kinetic-julia}

\begin{itemize}
  \item \texttt{ODEProblem(M, f, u0, tspan)} with Rosenbrock or SDIRK methods
  \item Supports mass matrix $M$ (constant or time-dependent)
  \item \textbf{Ref:} Rackauckas \& Nie (2017), JORS 5(1), 15
  \item \textbf{How it would work:} Export $K$ from Wolfram to JSON, construct $M$ in Julia, use \texttt{Rodas5()} or \texttt{Rosenbrock23()} for time integration
  \item \textbf{Advantage:} Handles stiff systems, implicit methods, arbitrary $M$
  \item \textbf{Disadvantage:} Requires Julia runtime, not Python; would need new backend
\end{itemize}

\subsubsection{PETSc TS}
\label{sec:kinetic-petsc}

\begin{itemize}
  \item \texttt{DMTSCreateRHSMassMatrix()} for non-identity mass matrix
  \item Implicit Runge--Kutta, BDF, ARKIMEX schemes
  \item \textbf{Ref:} Balay et al.\ (2024), PETSc/TAO Users Manual
  \item \textbf{Advantage:} Industrial-strength, scalable, MPI-parallel
  \item \textbf{Disadvantage:} Heavy dependency, C API, complex setup
\end{itemize}

\subsubsection{SUNDIALS IDA}
\label{sec:kinetic-ida}

\begin{itemize}
  \item General DAE solver: $F(t, y, y') = 0$ --- avoids $K^{-1}$ entirely
  \item Variable-order BDF with Newton iteration
  \item \textbf{Ref:} Hindmarsh et al.\ (2005), ACM TOMS 31(3)
  \item \textbf{Advantage:} Most general formulation, handles constraints naturally
  \item \textbf{Disadvantage:} Implicit (requires Jacobian), overkill for linear systems
\end{itemize}

\subsubsection{When to Consider}
\label{sec:kinetic-alt1-when}

\begin{itemize}
  \item Position-dependent $K$ (curved spacetime with non-constant metric)
  \item Very large $K$ ($>50\times 50$) where symbolic inversion produces huge expressions
  \item Stiff systems where explicit methods are CFL-limited
  \item If the numpy backend is replaced with a Julia or C++ backend
\end{itemize}


\subsection{Alternative 2: Gauge Fixing (Reduce $K$ to Diagonal)}
\label{sec:kinetic-alt2}

\textbf{Status:} TT gauge implemented (Phase B). De Donder implemented (Phase B). See \cref{sec:gauge-fixing}.

Instead of inverting $K$, apply gauge conditions that eliminate off-diagonal kinetic coupling, making $K$ diagonal (or even identity).

\subsubsection{TT Gauge (Transverse-Traceless)}
\label{sec:kinetic-tt}

\begin{itemize}
  \item Eliminates non-physical components entirely
  \item Reduces 10 GR components $\to$ 2 physical polarizations ($h_+$, $h_\times$)
  \item Type B constraint: temporal zeroing + traceless + transverse
  \item \textbf{Works for:} Gravitational waves in vacuum (massless)
  \item \textbf{Doesn't work for:} Massive gravity (need all components for mass term)
\end{itemize}

\subsubsection{De Donder (Harmonic) Gauge}
\label{sec:kinetic-dedonder}

\begin{itemize}
  \item Adds $\lag_\text{gf} = -\frac{1}{2\xi}\bigl(\partial^a h_{ab} - \tfrac{1}{2}\partial_b h\bigr)^2$ to the Lagrangian
  \item Diagonalizes $K$ and exposes constraint structure ($h_0$, $h_{0i}$ become constraints)
  \item Type A (Lagrangian term): implemented
  \item \textbf{Works for:} Both massless and massive gravity
  \item \textbf{Disadvantage:} Gauge-dependent; results must be checked for gauge artifacts
\end{itemize}

\subsubsection{When to Consider}
\label{sec:kinetic-alt2-when}

\begin{itemize}
  \item When physical interpretation of individual components is needed
  \item When reducing computational cost (fewer components to evolve)
  \item When connecting to observables in specific gauges (e.g., TT for GW detection)
\end{itemize}


\subsection{Alternative 3: Symplectic Integrators with Pre-Computed $K^{-1}$}
\label{sec:kinetic-alt3}

\textbf{Status:} Investigated, not implemented.

For constant $K$ (flat spacetime, constant parameters), pre-compute $K^{-1}$ once and use explicit symplectic integrator (leapfrog/Verlet) modified for $K^{-1}$:
\begin{align}
  q_{n+1}   &= q_n + \Delta t\; K^{-1} \cdot p_{n+1/2} \label{eq:symplectic-q} \\
  p_{n+1}   &= p_n + \Delta t\; F(q_{n+1}) \label{eq:symplectic-p}
\end{align}

\begin{itemize}
  \item \textbf{Ref:} Hairer \& Lubich (2003), ZAMM 83(1)
  \item \textbf{Advantage:} Energy-conserving (symplectic), explicit, no Jacobian
  \item \textbf{Disadvantage:} Only works for constant $K$; our current explicit RK already works
  \item \textbf{When to consider:} Long-time simulations where energy conservation is critical
\end{itemize}


\subsection{Alternative 4: Runtime Numeric Inversion (NumPy)}
\label{sec:kinetic-alt4}

\textbf{Status:} Not implemented. Simplest fallback.

Build $K$ numerically from JSON coefficients at runtime, invert with \texttt{numpy.linalg.inv}.

\begin{itemize}
  \item \textbf{Advantage:} Simple, no Wolfram dependency, handles parameter sweeps
  \item \textbf{Disadvantage:} Numerical (not exact), repeated inversion per timestep for position-dependent $K$, doesn't leverage symbolic structure
  \item \textbf{When to consider:} If symbolic $K^{-1}$ expressions become too large
\end{itemize}


\subsection{Alternative 5: Mass Lumping (FEM Approximation)}
\label{sec:kinetic-alt5}

\textbf{Status:} Not applicable to our spectral/FD discretization.

In FEM, the mass matrix $M$ arises from basis function integration. ``Mass lumping'' replaces $M$ with a diagonal approximation (row-sum or higher-order lumping).

\begin{itemize}
  \item \textbf{Ref:} Rathgeber et al.\ (2016), ACM TOMS 43(3)
  \item \textbf{Not applicable} because our $K$ comes from the Lagrangian structure, not from spatial discretization. $K$ is typically small ($10\times 10$ for GR) and exact inversion is feasible.
\end{itemize}


\subsection{Alternative 6: MultiModeCode-Style Field-Space Metric}
\label{sec:kinetic-alt6}

\textbf{Status:} Investigated for reference.

In multi-field inflation, the kinetic matrix is the ``field-space metric'' $G_{IJ}$. MultiModeCode evolves transport equations with non-trivial $G_{IJ}$ directly.

\begin{itemize}
  \item \textbf{Ref:} Dias et al.\ (2015), JCAP 2015(01):030
  \item \textbf{Relevance:} Similar mathematical structure to our $K_{ij}$
  \item \textbf{Difference:} They evolve ODEs (homogeneous fields), we evolve PDEs (spatial grid)
  \item \textbf{When to consider:} If we ever add multi-field inflation examples
\end{itemize}


\subsection{Decision Matrix}
\label{sec:kinetic-decision}

\begin{table}[htbp]
\centering
\caption{Comparison of kinetic matrix inversion approaches.}
\label{tab:kinetic-decision}
\begin{tabular}{@{}lccccc@{}}
\toprule
Approach & Implemented & Exact & Runtime Cost & Position-Dep.\ $K$ & Handles DAEs \\
\midrule
Symbolic $K^{-1}$ (current)      & Yes     & Yes & Zero   & Partial$^*$ & No  \\
Julia DifferentialEquations.jl   & No      & No  & Medium & Yes         & Yes \\
PETSc TS                         & No      & No  & Low    & Yes         & Yes \\
SUNDIALS IDA                     & No      & No  & Medium & Yes         & Yes \\
Gauge fixing                     & Partial & Yes & Zero   & N/A         & N/A \\
Symplectic $K^{-1}$              & No      & Yes & Low    & No          & No  \\
NumPy inversion                  & No      & No  & High   & Yes         & No  \\
\bottomrule
\end{tabular}
\end{table}

$^*$Position-dependent $K$: \texttt{coefficient\_symbolic} + \texttt{evaluate\_coefficient()} supports it, but not yet tested with non-constant $K^{-1}$ entries.
